\section{Beyond Comparing Measures}
\label{sec-beyond-comparing-measures}
\subsection{Vector and Matrix-Valued Measures}
\label{sec-vector-matrix-valued-measures}
\paragraph{Positive vector-valued measures.}
\begin{defn}[Positive vector-valued measure]\label{def-positive-vector-valued-measure}
A positive $\RR_+^m$-valued measure on $\X$ is a tuple
\(\al=(\al^1,\ldots,\al^m)\in\Mm_+(\X;\RR_+^m),\)
where each component $\al^k$ is a nonnegative finite measure.
\end{defn}
For the dynamic formulas below, assume that $\X$ is either $\RR^d$, the flat torus, or a bounded convex domain in $\RR^d$ equipped with no-flux boundary conditions. First suppose that the endpoints and the curve have densities.

\begin{equation}\label{eq-vector-valued-bb}
\mathcal W_{\Phi}^2(\al_0,\al_1)
\eqdef
\inf_{u,V}
\int_0^1\!\int_\X
\Phi(u_t(x),V_t(x))\,\d x\,\d t
\end{equation}
\begin{equation}\label{eq-vector-valued-continuity}
\partial_t u_t+\nabla_x\cdot V_t=0,
\qquad
(\nabla_x\cdot V_t)^k=\sum_{\ell=1}^d \partial_{x_\ell} V_{t,\ell}^k.
\end{equation}
\begin{prop}[Diagonal positive vector Benamou--Brenier]\label{prop-diagonal-positive-vector-bb}
Assume that $\al_0^k,\al_1^k\in\Mm_+(\X)$ have finite second moments and the same mass $m_k$ for every $k$. For the diagonal mobility $\mathsf M_{\mathrm{diag}}$, the value of~\eqref{eq-vector-valued-bb} is
\[
\mathcal W_{\mathrm{diag}}^2(\al_0,\al_1)
=
\sum_{k:m_k>0} m_k\,\Wass_2^2\!\left(\frac{\al_0^k}{m_k},\frac{\al_1^k}{m_k}\right),
\]
with the convention that zero-mass channels contribute zero.
\end{prop}
\begin{proof}
For $\mathsf M_{\mathrm{diag}}(u)=\diag(u_1,\ldots,u_m)$, the action separates as \(\sum_{\ell=1}^d V_\ell^\top \mathsf M_{\mathrm{diag}}(u)^\dagger V_\ell = \sum_{k=1}^m \frac{|V^k|^2}{u^k},\)
where $V^k=(V_1^k,\ldots,V_d^k)$ is the spatial momentum of channel $k$, and the scalar perspective convention is used. The constraint~\eqref{eq-vector-valued-continuity} also separates into $\partial_t u^k+\nabla\cdot V^k=0$. The minimization therefore splits into $m$ independent scalar Benamou--Brenier problems. If $m_k=0$, nonnegativity and conservation force the whole channel to vanish. If $m_k>0$, normalizing $\rho_t^k=u_t^k/m_k$ and $p_t^k=V_t^k/m_k$ factors the channel action as $m_k\int |p_t^k|^2/\rho_t^k$, hence the scalar value is $m_k\Wass_2^2(\al_0^k/m_k,\al_1^k/m_k)$. Summing over the channels proves the claim.
\end{proof}
\paragraph{Positive matrix-valued measures.}
\begin{defn}[Positive matrix-valued measure]\label{def-positive-matrix-valued-measure}
Write $\mathbb S^m$ for real symmetric matrices and $\mathbb S_+^m$ for the positive semidefinite cone. A positive $\mathbb S_+^m$-valued measure is a countably additive map from Borel sets of $\X$ to $\mathbb S_+^m$; we denote the cone of such measures by
\(\mathcal A\in\Mm_+(\X;\mathbb S_+^m).\)
\end{defn}
\begin{equation}\label{eq-matrix-valued-bb}
\begin{aligned}
\mathcal W_{\mathrm{mat}}^2(\mathcal A_0,\mathcal A_1)
\eqdef
\inf_{A,P}\int_0^1\!\int_\X
\sum_{\ell=1}^d \tr\big(P_{t,\ell}^{\top} A_t^\dagger P_{t,\ell}\big)\d x\d t
\end{aligned}
\end{equation}
\begin{equation}\label{eq-matrix-valued-continuity}
\partial_t A_t+\nabla_x\cdot P_t=0,
\qquad
\nabla_x\cdot P_t=\sum_{\ell=1}^d \partial_{x_\ell}P_{t,\ell}.
\end{equation}
\begin{prop}[Diagonal matrix subproblem]\label{prop-matrix-diagonal-reduction}
Assume that the endpoints are diagonal in a fixed orthonormal basis, \(\mathcal A_i=\diag(\al_i^1,\ldots,\al_i^m), \qquad i=0,1,\)
and that $\al_0^k(\X)=\al_1^k(\X)=m_k$ for every $k$. If one restricts the admissible curves in~\eqref{eq-matrix-valued-bb} to remain diagonal in that basis,
\(A_t=\diag(u_t^1,\ldots,u_t^m), \qquad P_{t,\ell}=\diag(V_{t,\ell}^1,\ldots,V_{t,\ell}^m),\)
then the value of this restricted matrix problem is
\[
\sum_{k:m_k>0} m_k\,\Wass_2^2\!\left(\frac{\al_0^k}{m_k},\frac{\al_1^k}{m_k}\right),
\]
with zero contribution from zero-mass channels. Thus the commuting matrix submodel is exactly the diagonal positive vector-valued Benamou--Brenier model of Proposition~\ref{prop-diagonal-positive-vector-bb}.
\end{prop}
\begin{proof}
The continuity equation~\eqref{eq-matrix-valued-continuity} is diagonal entry by diagonal entry and gives $\partial_t u^k+\nabla\cdot V^k=0$. Moreover,
\(\sum_{\ell=1}^d \tr\big(P_{t,\ell}^{\top}A_t^\dagger P_{t,\ell}\big) = \sum_{k=1}^m \frac{|V_t^k|^2}{u_t^k}\)
with the same scalar perspective convention as before. The admissible set and the action are therefore exactly those of the diagonal vector model.
\end{proof}
\subsection{Wasserstein over Wasserstein}
\label{sec-wasserstein-over-wasserstein}
\begin{prop}[Wasserstein spaces as ground spaces]\label{prop-wasserstein-space-polish}
Let $1\leq p<\infty$. If $(\X,\dist)$ is Polish, then $(\Pp_p(\X),\Wass_p)$ is Polish. If $\X$ is compact, then $\Pp_p(\X)=\Pp(\X)$ and this space is compact for $\Wass_p$.
\end{prop}
\begin{proof}
This is a standard structural theorem for Wasserstein spaces. For completeness, extract from a $\Wass_p$-Cauchy sequence a subsequence $(\alpha_k)_k$ such that $\sum_k\Wass_p(\alpha_k,\alpha_{k+1})<\infty$. Gluing optimal couplings of consecutive terms produces random variables $(X_k)_k$ with laws $\alpha_k$ and $\sum_k\norm{\dist(X_k,X_{k+1})}_{L^p}<\infty$. Hence $X_k$ converges almost surely and in $L^p$ to an $\X$-valued random variable; its law is the $\Wass_p$ limit of the subsequence, and the Cauchy property gives convergence of the original sequence. Separability follows by approximating measures with finitely supported measures on a countable dense subset and rational weights. If $\X$ is compact, Prokhorov's theorem gives compactness of $\Pp(\X)$ for weak convergence, and Proposition~\ref{prop-wass-metrizes-weak-compact} identifies this topology with the $\Wass_p$ topology.
\end{proof}
Fix $1\leq p<\infty$. We denote elements of $\Pp_p(\X)$ by $\alpha,\beta,\ldots$. Elements of $\Pp_p(\Pp_p(\X))$ are denoted by fraktur letters, for instance $\mathfrak A,\mathfrak B$; they are probability laws over probability measures, or random probability measures. A basic parametric example is obtained from a measurable family $(\alpha_\zeta)_{\zeta\in Z}$ and a probability law $\gamma$ on the parameter space:

\begin{equation}\label{eq-wow-parametric-law}
\mathfrak A=(\zeta\mapsto\alpha_\zeta)_\sharp\gamma.
\end{equation}
\(\int_Z \Wass_p(\alpha_\zeta,\delta_{x_0})^p\,\d\gamma(\zeta)<\infty.\)
If $\gamma=\sum_{i=1}^K a_i\delta_{\zeta_i}$, then

\(\mathfrak A=\sum_{i=1}^K a_i\delta_{\alpha_{\zeta_i}}.\)
\begin{defn}[Collapsed, or barycentric, mixture]\label{def-collapsed-barycentric-mixture}
For $\mathfrak A\in\Pp(\Pp(\X))$, the collapsed, or barycentric, mixture associated with $\mathfrak A$ is the measure $\bar\alpha_{\mathfrak A}$ defined by
\begin{equation}\label{eq-wow-barycentric-mixture}
\int_\X f(x)\d\bar\alpha_{\mathfrak A}(x)
=
\int_{\Pp(\X)}
\left(\int_\X f(x)\d\alpha(x)\right)
\d\mathfrak A(\alpha),
\end{equation}
for bounded continuous $f$.
\end{defn}
\(\int_\X \dist(x,x_0)^p\,\d\bar\alpha_{\mathfrak A}(x) = \int_{\Pp_p(\X)}\Wass_p(\alpha,\delta_{x_0})^p\,\d\mathfrak A(\alpha).\)
\begin{equation}\label{eq-wow-distance}
\mathbb W_p^p(\mathfrak A,\mathfrak B)
\eqdef
\inf_{\Pi\in\Couplings(\mathfrak A,\mathfrak B)}
\int_{\Pp_p(\X)\times\Pp_p(\X)}
\Wass_p^p(\alpha,\beta)\d\Pi(\alpha,\beta),
\qquad
\mathfrak A,\mathfrak B\in\Pp_p(\Pp_p(\X)).
\end{equation}
For $p=2$, Gaussian mixtures provide a particularly explicit example with two levels of geometry. A mixture $\sum_i a_i\Gaussian(m_i,\Sigma_i)$ can either be viewed as the collapsed measure on $\X$, or as the component law

\(\mathfrak A=\sum_i a_i\delta_{\Gaussian(m_i,\Sigma_i)}\)
\(\mathfrak A=\sum_i a_i\delta_{\Gaussian(m_i,\Sigma_i)}, \qquad \mathfrak B=\sum_j b_j\delta_{\Gaussian(n_j,\Lambda_j)},\)
\(C_{ij}=\norm{m_i-n_j}^2+\Bb(\Sigma_i,\Lambda_j)^2.\)
Let $\Pi^\star$ be an optimal coupling between the weights $a$ and $b$. If $A_{ij}$ denotes the Brenier linear part from $\Sigma_i$ to $\Lambda_j$, then each active component pair is interpolated by

\(m_{ij,t}=(1-t)m_i+t n_j, \qquad \Sigma_{ij,t} = \big((1-t)\Id+tA_{ij}\big)\Sigma_i \big((1-t)\Id+tA_{ij}\big)^\top,\)
\(\bar\alpha_t= \sum_{i,j}\Pi^\star_{ij}\Gaussian(m_{ij,t},\Sigma_{ij,t}).\)
\begin{prop}[Collapsing is non-expansive]\label{prop-wow-collapsed-bound}
Let $(\X,\dist)$ be Polish, let $1\leq p<\infty$, and let $\mathfrak A,\mathfrak B\in\Pp_p(\Pp_p(\X))$. For the collapsed mixtures defined by~\eqref{eq-wow-barycentric-mixture},
\(\Wass_p(\bar\alpha_{\mathfrak A},\bar\beta_{\mathfrak B}) \leq \mathbb W_p(\mathfrak A,\mathfrak B).\)
\end{prop}
\begin{proof}
Fix $\Pi\in\Couplings(\mathfrak A,\mathfrak B)$ and $\eta>0$. A standard measurable-selection theorem for optimal transport on Polish spaces gives a measurable kernel $(\alpha,\beta)\mapsto\pi_{\alpha,\beta}\in\Couplings(\alpha,\beta)$ such that
\(\int_{\X\times\X}\dist(x,y)^p\,\d\pi_{\alpha,\beta}(x,y) \leq \Wass_p(\alpha,\beta)^p+\eta.\)
Integrating this kernel against $\Pi$ gives a coupling $\bar\pi$ between $\bar\alpha_{\mathfrak A}$ and $\bar\beta_{\mathfrak B}$. Its cost satisfies
\(\int_{\X\times\X}\dist(x,y)^p\d\bar\pi(x,y) \leq \int_{\Pp_p(\X)^2}\Wass_p^p(\alpha,\beta)\d\Pi(\alpha,\beta)+\eta.\)
Letting $\eta\downarrow0$, taking the infimum over $\Pi$, and then taking the $p$-th root proves the claim.
\end{proof}
\subsection{Gromov--Wasserstein}
\label{sec-gromov-wasserstein}
\paragraph{Discrete formulation.}
Optimal transport needs a ground cost $\C$ to compare histograms $(\a,\b)$, and thus cannot be used directly if the histograms are not defined on the same underlying space, or if one cannot pre-register these spaces to define a ground cost.

To address this issue, one can instead use a weaker requirement: two matrices $\distD \in \RR^{n \times n}$ and $\distD' \in \RR^{m \times m}$ are available and represent relationships between the points on which the histograms are defined. A typical scenario is when these matrices are powers of distance matrices.

\begin{equation}\label{eq-gw-def}
\begin{aligned}
\Ee_{\distD,\distD'}(\P)
&\eqdef
\sum_{i,j,i',j'} \De(\distD_{i,i'},\distD'_{j,j'})^p \P_{i,j}\P_{i',j'},
\\
\GWD( (\a,\distD), (\b,\distD') )^p
&\eqdef
\min_{\P \in \CouplingsD(\a,\b) }
\Ee_{\distD,\distD'}(\P),
\end{aligned}
\end{equation}
where $p \geq 1$ and $\De$ is a distance on $\RR$, typically $\De(u,v)=|u-v|$.

When the matrices $\distD,\distD'$ are genuine distance matrices, the general construction below shows that $\GWD$ satisfies the triangle inequality and defines a distance between metric spaces equipped with a probability distribution, up to measure-preserving isometries.

\paragraph{General setting.}
\begin{defn}[Metric-measure space]\label{def-metric-measure-space}
A metric-measure space is a triple
\(\XX=(\X,\dist_\X,\al),\)
where $(\X,\dist_\X)$ is Polish and $\al$ is a Borel probability measure on $\X$. For a finite $p$-Gromov--Wasserstein distance, one additionally assumes the finite-size condition
\(\int_{\X\times\X}\dist_\X(x,x')^p\,\d\al(x)\d\al(x')<\infty.\)
\end{defn}
The general setting corresponds to computing couplings between metric-measure spaces $\XX=(\X,\dist_\X,\al)$ and $\YY=(\Y,\dist_\Y,\be)$, where the distance and the measure are both part of the data. Compactness is often assumed below to avoid additional tightness and integrability arguments.

\begin{align}
\label{eq-gw-generic}
\GW( \XX,\YY )^p \eqdef
\umin{ \pi \in \Couplings(\al,\be) }
\int_{\X^2 \times \Y^2}
\De(\dist_\X(x,x'),\dist_\Y(y,y'))^p
\d\pi(x,y)\d\pi(x',y').
\end{align}
\begin{prop}[Fixed-space GW is controlled by Wasserstein]\label{prop-gw-controlled-by-wasserstein}
Let $\alpha,\beta\in\Pp_p(\X)$ be probability measures on the same metric space $(\X,\dist_\X)$, and take $\De(u,v)=|u-v|$ in~\eqref{eq-gw-generic}. Then
\(\GW((\X,\dist_\X,\alpha),(\X,\dist_\X,\beta)) \leq 2\Wass_p(\alpha,\beta).\)
The Wasserstein distance on the right is computed with the same ground metric $\dist_\X$. The factor $2$ comes from the normalization in~\eqref{eq-gw-generic}; with the alternative convention $\De(u,v)=|u-v|/2$, the displayed constant becomes $1$.
\end{prop}
\begin{proof}
Let $\pi$ be any coupling between $\alpha$ and $\beta$. For two independent pairs $(X,Y),(X',Y')\sim\pi$, the reverse triangle inequality gives \(\big|\dist_\X(X,X')-\dist_\X(Y,Y')\big| \leq \dist_\X(X,Y)+\dist_\X(X',Y').\)
Taking the $L^p$ norm and using Minkowski gives a bound by $2(\int\dist_\X(x,y)^p\d\pi)^{1/p}$. Optimizing over $\pi$ proves the claim.
\end{proof}
\begin{defn}[Isometric metric-measure spaces]\label{def-isometric-mm-spaces}
Two metric-measure spaces $\XX=(\X,\dist_\X,\al)$ and $\YY=(\Y,\dist_\Y,\be)$ are isometric if there exists a bijection $\phi:\supp(\al)\to\supp(\be)$ such that $\phi_\sharp\al=\be$ and
\(\dist_\Y(\phi(x),\phi(x'))=\dist_\X(x,x')\)
for all $x,x'\in\supp(\al)$.
\end{defn}
\begin{thm}[Gromov--Wasserstein metric modulo isometries]\label{thm-gw-metric}
For compact metric-measure spaces, $p\geq1$ and $\De(u,v)=|u-v|$, $\GW$ defines a distance up to measure-preserving isometries.
\end{thm}
\begin{proof}
Compactness of the supports makes the coupling set compact and the distortion functional continuous, so an optimal coupling exists. If $\phi$ is a measure-preserving isometry, the graph coupling $(\Id,\phi)_\sharp\al$ has zero distortion; in particular, $\GW(\XX,\XX)=0$. Symmetry and non-negativity are immediate.

If $\GW(\XX,\YY)=0$ and $\pi$ is an optimal plan, then $\dist_\X(x,x')=\dist_\Y(y,y')$ holds $\pi\otimes\pi$-almost everywhere. By continuity, this equality holds on $\supp(\pi)^2$.
\(\dist_\pi((x,y),(x',y')) \eqdef \frac12\dist_\X(x,x')+\frac12\dist_\Y(y,y').\)
The first projection $\psi:(x,y)\mapsto x$ is measure-preserving. For $((x,y),(x',y'))\in\supp(\pi)^2$, \(\dist_\X(\psi(x,y),\psi(x',y')) = \dist_\X(x,x') = \dist_\Y(y,y') = \dist_\pi((x,y),(x',y')),\)
so $\psi$ is an isometry and therefore injective. To see surjectivity onto $\supp(\al)$, take $x\in\supp(\al)$. Since $\psi_\sharp\pi=\al$, there is a sequence $(x_k,y_k)\in\supp(\pi)$ with $x_k\to x$. The equality of distances on $\supp(\pi)$ makes $(y_k)_k$ Cauchy, and compactness gives a convergent subsequence with limit $(x,y)\in\supp(\pi)$. The same argument for the second projection shows that the support space is also isometric to $\YY$.

For the triangle inequality, let $\pi$ be an optimal coupling between $\XX$ and $\YY$, and $\xi$ an optimal coupling between $\YY$ and $\ZZ=(\Z,\dist_\Z,\ga)$. By the gluing lemma, take $\sigma$ on $\X\times\Y\times\Z$ whose $(\X,\Y)$ and $(\Y,\Z)$ marginals are $\pi$ and $\xi$. Let $\rho=(P_{\X,\Z})_\sharp\sigma$, and write $\bar\sigma=\sigma\otimes\sigma$ for the product law of two independent triples $(x,y,z)$ and $(x',y',z')$. Then $\rho$ is feasible between $\XX$ and $\ZZ$, and
\begin{align*}
\GW(\XX,\ZZ)
&\leq
\left(
\int
\abs{\dist_\X(x,x')-\dist_\Z(z,z')}^p
\d\bar\sigma
\right)^{1/p} \\
&\leq
\left(
\int
\abs{\dist_\X(x,x')-\dist_\Y(y,y')}^p
\d\bar\sigma
\right)^{1/p}
+
\left(
\int
\abs{\dist_\Y(y,y')-\dist_\Z(z,z')}^p
\d\bar\sigma
\right)^{1/p} \\
&=
\GW(\XX,\YY)+\GW(\YY,\ZZ),
\end{align*}
where the second inequality uses the pointwise triangle inequality followed by Minkowski's inequality. This proves the triangle inequality and completes the metric proof.
\end{proof}
\begin{prop}[Marginal stability and empirical GW rates]\label{prop-gw-empirical-stability}
Let $\XX=(\X,\dist_\X,\alpha)$ and $\YY=(\Y,\dist_\Y,\beta)$ be compact metric-measure spaces, with $\De(u,v)=|u-v|$ and exponent $p\geq1$. For any probability measures $\tilde\alpha$ on $\X$ and $\tilde\beta$ on $\Y$, set
\(\widetilde\XX=(\X,\dist_\X,\tilde\alpha), \qquad \widetilde\YY=(\Y,\dist_\Y,\tilde\beta).\)
Then, deterministically, \(\abs{\GW(\widetilde\XX,\widetilde\YY)-\GW(\XX,\YY)} \leq 2\Wass_p^\X(\tilde\alpha,\alpha) + 2\Wass_p^\Y(\tilde\beta,\beta),\)
where $\Wass_p^\X$ and $\Wass_p^\Y$ denote Wasserstein distances computed with $\dist_\X$ and $\dist_\Y$. If $\tilde\alpha=\hat\alpha_n$ and $\tilde\beta=\hat\beta_m$ are independent empirical measures built from iid samples with laws $\alpha$ and $\beta$, and if
\(\widehat\XX_n=(\X,\dist_\X,\hat\alpha_n), \qquad \widehat\YY_m=(\Y,\dist_\Y,\hat\beta_m),\)
then
\(\EE\abs{\GW(\widehat\XX_n,\widehat\YY_m)-\GW(\XX,\YY)} \leq 2\EE\Wass_p^\X(\hat\alpha_n,\alpha) + 2\EE\Wass_p^\Y(\hat\beta_m,\beta).\)
In particular, if $\X$ and $\Y$ are regular compact subsets of Euclidean spaces of intrinsic dimensions $d_\X$ and $d_\Y$, and the high-dimensional Wasserstein regime $d_\X,d_\Y>2p$ applies, then
\(\EE\abs{\GW(\widehat\XX_n,\widehat\YY_m)-\GW(\XX,\YY)} = O(n^{-1/d_\X}+m^{-1/d_\Y}).\)
For $p=1$ on subsets of $[0,1]^d$, the rates, including the low-dimensional logarithmic cases, are those of Proposition~\ref{prop-empirical-ot-rate}.
\end{prop}
\begin{proof}
The triangle inequality of Theorem~\ref{thm-gw-metric} gives \(\GW(\widetilde\XX,\widetilde\YY) - \GW(\XX,\YY) \leq \GW(\widetilde\XX,\XX) + \GW(\YY,\widetilde\YY).\)
Exchanging the two pairs gives the same bound for the opposite difference, hence the absolute value. The two terms on the right compare metric-measure spaces with the same underlying metric and only different measures, so Proposition~\ref{prop-gw-controlled-by-wasserstein} bounds them by $2\Wass_p^\X(\tilde\alpha,\alpha)$ and $2\Wass_p^\Y(\tilde\beta,\beta)$. Substituting $\tilde\alpha=\hat\alpha_n$ and $\tilde\beta=\hat\beta_m$ and taking expectations gives the second display. The Euclidean rate then follows from the empirical Wasserstein rates recalled in Section~\ref{sec-sample-complexity}.
\end{proof}
\begin{prop}[Gromov--Wasserstein geodesics]\label{prop-gw-geodesics}
Let $\XX_0=(\X_0,\dist_{\X_0},\al_0)$ and $\XX_1=(\X_1,\dist_{\X_1},\al_1)$ be compact metric-measure spaces, take $\De(u,v)=|u-v|$ in~\eqref{eq-gw-generic}, and let $\pi^\star$ be an optimal coupling. Define, on $\Z=\X_0\times\X_1$,
\(\dist_t((x_0,x_1),(x'_0,x'_1)) \eqdef (1-t)\dist_{\X_0}(x_0,x'_0)+t\dist_{\X_1}(x_1,x'_1), \qquad \XX_t=(\Z,\dist_t,\pi^\star).\)
For $0<t<1$, $\dist_t$ is a metric. At $t=0$ and $t=1$, one quotients $\Z$ by the zero-distance relation associated with $\dist_t$. Then $t\mapsto\XX_t$ is a constant-speed geodesic:
\(\GW(\XX_s,\XX_t)=|t-s|\GW(\XX_0,\XX_1) \qquad \forall s,t\in[0,1].\)
\end{prop}
\begin{proof}
Write $D=\GW(\XX_0,\XX_1)$. For $s<t$, couple $\XX_s$ and $\XX_t$ by the diagonal coupling induced by the identity on $\Z$ and the measure $\pi^\star$. For two independent points $z=(x_0,x_1)$ and $z'=(x'_0,x'_1)$ sampled from $\pi^\star$,
\(\dist_t(z,z')-\dist_s(z,z') = (t-s)\big(\dist_{\X_1}(x_1,x'_1)-\dist_{\X_0}(x_0,x'_0)\big).\)
Using this feasible coupling gives $\GW(\XX_s,\XX_t)\leq(t-s)D$. The same construction with the projections from $\Z$ to $\X_0$ and $\X_1$ gives $\GW(\XX_0,\XX_t)\leq tD$ and $\GW(\XX_t,\XX_1)\leq(1-t)D$. The triangle inequality for $\GW$ then yields, for $0\leq s\leq t\leq1$,
\(D \leq \GW(\XX_0,\XX_s)+\GW(\XX_s,\XX_t)+\GW(\XX_t,\XX_1) \leq sD+\GW(\XX_s,\XX_t)+(1-t)D,\)
so $\GW(\XX_s,\XX_t)\geq(t-s)D$. This proves equality.
\end{proof}
\begin{prop}[M\'emoli profile lower bound]\label{prop-memoli-gw-profile-lower-bound}
Let $\XX=(\X,\dist_\X,\alpha)$ and $\YY=(\Y,\dist_\Y,\beta)$ be compact metric-measure spaces and take $\De(u,v)=|u-v|$ in~\eqref{eq-gw-generic}, with the same exponent $p\geq1$. For each $x\in\X$ and $y\in\Y$, define the distance-profile measures on $\RR_+$ by
\(\alpha_x\eqdef(\dist_\X(x,\cdot))_\sharp\alpha, \qquad \beta_y\eqdef(\dist_\Y(y,\cdot))_\sharp\beta.\)
Let $\mathfrak D_\XX=(x\mapsto\alpha_x)_\sharp\alpha$ and $\mathfrak D_\YY=(y\mapsto\beta_y)_\sharp\beta$, which are probability measures on $\Pp_p(\RR_+)$. Then
\(\mathbb W_p\bigl(\mathfrak D_\XX,\mathfrak D_\YY\bigr) \leq \GW(\XX,\YY).\)
Here $\mathbb W_p$ is the Wasserstein-over-Wasserstein distance~\eqref{eq-wow-distance}: its ground space is $(\Pp_p(\RR_+),\Wass_p)$.
\end{prop}
\begin{proof}
Fix any $\pi\in\Couplings(\alpha,\beta)$. It induces a coupling $(x,y)\mapsto(\alpha_x,\beta_y)$ between $\mathfrak D_\XX$ and $\mathfrak D_\YY$, hence
\(\mathbb W_p\bigl(\mathfrak D_\XX,\mathfrak D_\YY\bigr)^p \leq \int_{\X\times\Y}\Wass_p(\alpha_x,\beta_y)^p\,\d\pi(x,y).\)
For fixed $(x,y)$, the map $(x',y')\mapsto(\dist_\X(x,x'),\dist_\Y(y,y'))$ pushes the same coupling $\pi$ to a coupling between $\alpha_x$ and $\beta_y$. Therefore
\(\Wass_p(\alpha_x,\beta_y)^p \leq \int_{\X\times\Y} \abs{\dist_\X(x,x')-\dist_\Y(y,y')}^p \d\pi(x',y').\)
Integrating in $(x,y)$ gives \(\mathbb W_p\bigl(\mathfrak D_\XX,\mathfrak D_\YY\bigr)^p \leq \int_{\X^2\times\Y^2} \abs{\dist_\X(x,x')-\dist_\Y(y,y')}^p \d\pi(x,y)\d\pi(x',y').\)
Taking the infimum over $\pi$ and then the $p$-th root proves the claim.
\end{proof}
\paragraph{Relation with Wasserstein--Procrustes.}
\begin{prop}[Wasserstein--Procrustes upper certificate]\label{prop-gw-procrustes-upper-certificate}
Let $\alpha,\beta$ be probability measures on $\RR^d$, equipped with the Euclidean distance, and take $\De(u,v)=|u-v|$ in~\eqref{eq-gw-generic}. If $\Wass_{p,\mathrm E(d)}$ denotes the quotient Wasserstein distance of Definition~\ref{def-quotient-wasserstein} for the Euclidean group, then
\begin{equation}\label{eq-gw-procrustes-upper-gw-section}
\GW((\RR^d,\norm{\cdot},\alpha),(\RR^d,\norm{\cdot},\beta))
\leq
2\,\Wass_{p,\mathrm E(d)}([\alpha],[\beta]).
\end{equation}
\end{prop}
\begin{proof}
Let $g,h\in\mathrm E(d)$ be rigid motions. They preserve all pairwise Euclidean distances, so
\((\RR^d,\norm{\cdot},\alpha) \quad\text{is isometric to}\quad (\RR^d,\norm{\cdot},g_\sharp\alpha),\)
and similarly for $\beta$ and $h_\sharp\beta$. Hence $\GW$ is unchanged by pushing the two measures forward by $g$ and $h$. Applying Proposition~\ref{prop-gw-controlled-by-wasserstein} to $g_\sharp\alpha$ and $h_\sharp\beta$ gives
\(\GW((\RR^d,\norm{\cdot},\alpha),(\RR^d,\norm{\cdot},\beta)) \leq 2\,\Wass_p(g_\sharp\alpha,h_\sharp\beta).\)
Taking the infimum over $g,h\in\mathrm E(d)$ gives~\eqref{eq-gw-procrustes-upper-gw-section}.
\end{proof}
\begin{equation}\label{eq-gw-profile-procrustes-sandwich}
\mathbb W_p(\mathfrak D_\XX,\mathfrak D_\YY)
\leq
\GW(\XX,\YY)
\leq
2\,\Wass_{p,\mathrm E(d)}([\alpha],[\beta]),
\end{equation}
where $\XX=(\RR^d,\norm{\cdot},\alpha)$ and $\YY=(\RR^d,\norm{\cdot},\beta)$. The left term is intrinsic and inexpensive, while the right term is an extrinsic rigid-registration certificate.

\paragraph{Entropic regularization and iterative solver.}
For the common squared distortion $\De(u,v)^2=(u-v)^2$, one often seeks a stationary point of the entropic relaxation

\eql{\label{eq-gw-entropy}
\umin{\P\in\CouplingsD(\a,\b)}
\Ee_{\distD,\distD'}(\P)-\epsilon\HD(\P).
}
\eql{\label{eq-gw-sinkh}
\C(\P)
\eqdef
\distD^{\odot2}\a\,\ones_m^\top
+
\ones_n(\distD'^{\odot2}\b)^\top
-
2\distD\,\P\,\transp{\distD'},
\qquad
\nabla\Ee_{\distD,\distD'}(\P)=2\C(\P).
}
\(\P^{(k+1)} = \uargmin{\P\in\CouplingsD(\a,\b)} \dotp{\C(\P^{(k)})}{\P}-\frac{\epsilon}{2}\HD(\P).\)
The factor $\epsilon/2$ is essential because~\eqref{eq-gw-sinkh} is one half of the gradient of the quadratic term. Each update is an ordinary entropic OT problem and can therefore be solved with Sinkhorn iterations.

\paragraph{Adding features.}
\[
\operatorname{FGW}_{\lambda,p}((\a,\distD),(\b,\distD'))^p
\eqdef
\umin{\P\in\CouplingsD(\a,\b)}
(1-\lambda)\sum_{i,j}M_{ij}\P_{ij}
+\lambda
\sum_{i,j,i',j'}\De(\distD_{ii'},\distD'_{jj'})^p\P_{ij}\P_{i'j'}.
\]
\paragraph{Hausdorff and Gromov--Hausdorff viewpoints.}
If $A,B$ are compact subsets of a common metric space $(\Z,\dist_\Z)$, their Hausdorff distance is

\[
d_{\mathrm H}^{\Z}(A,B)
=
\max\left\{
\sup_{a\in A}\inf_{b\in B}\dist_\Z(a,b),
\sup_{b\in B}\inf_{a\in A}\dist_\Z(a,b)
\right\}.
\]
\(d_{\mathrm{GH}}(\X,\Y) = \inf_{\Z,\phi,\psi} d_{\mathrm H}^{\Z}(\phi(\X),\psi(\Y)).\)
\subsection{Quantum Optimal Transport}
\label{sec-quantum-ot}
\paragraph{Finite-dimensional states and couplings.}
\begin{defn}[Hermitian and density matrices]\label{def-hermitian-density-matrices}
Let $\mathbb H_n$ be the real vector space of $n\times n$ Hermitian matrices, \(\mathbb H_n^+ = \enscond{A\in\mathbb H_n}{A\succeq0}, \qquad \mathbb H_n^{+,1} = \enscond{A\in\mathbb H_n^+}{\tr(A)=1}.\)
Elements of $\mathbb H_n^{+,1}$ are density matrices.
\end{defn}
A joint quantum state between $\CC^n$ and $\CC^m$ is a matrix $T\in\mathbb H_{nm}^+$ acting on $\CC^n\otimes\CC^m$. Its marginals are the partial traces, defined by duality through

\begin{equation}\label{eq-qot-partial-traces}
\tr(F\,\operatorname{Tr}_B T)=\tr((F\otimes\Id_m)T),
\qquad
\tr(G\,\operatorname{Tr}_A T)=\tr((\Id_n\otimes G)T),
\end{equation}
for all $F\in\mathbb H_n$ and $G\in\mathbb H_m$. Thus $\operatorname{Tr}_B(T)\in\mathbb H_n^+$ and $\operatorname{Tr}_A(T)\in\mathbb H_m^+$ play exactly the role of the two marginals of a classical coupling.

\begin{defn}[Finite-dimensional quantum OT]\label{def-finite-dimensional-qot}
Let $A\in\mathbb H_n^{+,1}$, $B\in\mathbb H_m^{+,1}$ and let $C\in\mathbb H_{nm}$ be a Hermitian cost observable. The quantum OT value is the semidefinite program
\begin{equation}\label{eq-qot-primal}
\mathrm{QOT}_C(A,B)
\eqdef
\min_{T\in\mathbb H_{nm}^+}
\enscond{\tr(CT)}{\operatorname{Tr}_B(T)=A,\ \operatorname{Tr}_A(T)=B}.
\end{equation}
\end{defn}
\begin{prop}[Quantum Kantorovich duality]\label{prop-qot-duality}
For $A\in\mathbb H_n^{+,1}$ and $B\in\mathbb H_m^{+,1}$, the dual of~\eqref{eq-qot-primal} is
\begin{equation}\label{eq-qot-dual}
\mathrm{QOT}_C(A,B)
=
\sup_{F\in\mathbb H_n,\ G\in\mathbb H_m}
\left\{
\tr(FA)+\tr(GB)
:\,
F\otimes\Id_m+\Id_n\otimes G\preceq C
\right\}.
\end{equation}
There is no duality gap. If $A$ and $B$ are positive definite, the supremum is attained. For singular marginals, the primal reduces to the tensor product of their supports; on this reduced space the corresponding dual maximum is attained, while the full-space formula above is always valid as a supremum.
\end{prop}
\begin{proof}
Introduce Hermitian Lagrange multipliers $F$ and $G$ for the two marginal constraints. The Lagrangian is
\[
\tr(CT)+\tr\!\left(F(A-\operatorname{Tr}_B T)\right)
+\tr\!\left(G(B-\operatorname{Tr}_A T)\right)
=
\tr(FA)+\tr(GB)
+\tr\!\left((C-F\otimes\Id_m-\Id_n\otimes G)T\right),
\]
where~\eqref{eq-qot-partial-traces} was used in the last equality. Minimizing over $T\succeq0$ gives a finite lower bound if and only if $C-F\otimes\Id_m-\Id_n\otimes G\succeq0$, in which case the infimum in $T$ is $0$. When $A,B\succ0$, the coupling $A\otimes B$ is strictly feasible, so Slater's theorem gives equality of primal and dual values and dual attainment.

For singular marginals, let $P_A$ and $P_B$ be their support projections. If $T$ is feasible, then
\(\tr\bigl(((\Id_n-P_A)\otimes\Id_m)T\bigr) = \tr((\Id_n-P_A)A)=0.\)
Positivity implies that $T$ vanishes on $(\supp A)^\perp\otimes\CC^m$; the analogous argument for $B$ yields
\(T=(P_A\otimes P_B)T(P_A\otimes P_B).\)
Thus the primal is exactly the problem obtained by compressing $A$, $B$ and $C$ to $\supp A\otimes\supp B$. The reduced marginals are positive definite, so Slater applies there. Equivalently, the unreduced dual values approach the reduced maximum after assigning sufficiently negative potentials on the orthogonal complements. This proves the stated supremum formula without a gap.
\end{proof}
\paragraph{Entropic regularization and Bregman iterations.}
\begin{defn}[von Neumann quantum entropy]\label{def-von-neumann-quantum-entropy}
For a density matrix or positive semidefinite matrix $T$, the shifted von Neumann entropy functional used here is
\[
H(T)=\tr\!\left(T(\log T-\Id)\right),
\qquad
\nabla H(T)=\log T,
\]
with the convention $0\log0=0$ on eigenvalues. This is the convex negative quantum entropy; on trace-one states it differs from the physical entropy $-\tr(T\log T)$ by a sign and an additive constant.
\end{defn}
For $\epsilon>0$ define

\begin{equation}\label{eq-qot-entropic-primal}
\mathrm{QOT}_C^\epsilon(A,B)
=
\min_{T\succeq0}
\left\{
\tr(CT)+\epsilon H(T)
:\,
\operatorname{Tr}_B(T)=A,
\operatorname{Tr}_A(T)=B
\right\}.
\end{equation}
\begin{prop}[Entropic quantum OT duality]\label{prop-qot-entropic-duality}
Assume $A\succ0$, $B\succ0$ and $\epsilon>0$. Then~\eqref{eq-qot-entropic-primal} has a unique positive minimizer. Its dual is
\begin{equation}\label{eq-qot-entropic-dual}
\mathrm{QOT}_C^\epsilon(A,B)
=
\max_{F\in\mathbb H_n,\ G\in\mathbb H_m}
\left\{
\tr(FA)+\tr(GB)
-\epsilon\,
\tr\exp\!\left(\frac{F\otimes\Id_m+\Id_n\otimes G-C}{\epsilon}\right)
\right\}.
\end{equation}
At optimality, primal and dual variables are linked by the Gibbs formula
\begin{equation}\label{eq-qot-gibbs-coupling}
T_e(F,G)
=
\exp\!\left(\frac{F\otimes\Id_m+\Id_n\otimes G-C}{\epsilon}\right),
\end{equation}
with $\operatorname{Tr}_B(T_e)=A$ and $\operatorname{Tr}_A(T_e)=B$.
\end{prop}
\begin{proof}
The feasible set is compact and nonempty, and it contains the positive definite point $A\otimes B$. The trace entropy is strictly convex on positive semidefinite matrices, hence the regularized primal has a unique minimizer. This minimizer is positive definite: if it had a non-trivial kernel, moving toward $A\otimes B$ would have entropy directional derivative $-\infty$ at the boundary while changing the linear cost only at finite rate. Slater's condition justifies the Lagrange dual computation. The Fenchel identity
\(\sup_{T\succeq0}\ \tr(YT)-\epsilon H(T) = \epsilon\,\tr\exp(Y/\epsilon)\)
is the matrix analogue of the scalar exponential conjugacy. Applying it to the Lagrangian of~\eqref{eq-qot-entropic-primal}, with
$Y=F\otimes\Id_m+\Id_n\otimes G-C$, gives~\eqref{eq-qot-entropic-dual}. The stationarity condition of this Fenchel identity gives~\eqref{eq-qot-gibbs-coupling}; differentiating the dual objective with respect to $F$ and $G$ yields the two marginal equations.
\end{proof}
\[
D_H(T|K)
=
\tr\!\left(T(\log T-\log K)-T+K\right).
\]
\(\mathcal M_A=\{T\succeq0:\operatorname{Tr}_B(T)=A\}, \qquad \mathcal M_B=\{T\succeq0:\operatorname{Tr}_A(T)=B\}.\)
\begin{prop}[Exact Bregman projections]\label{prop-qot-bregman-projections}
Assume $A,B\succ0$ and let $K=\exp(-C/\epsilon)$. The minimizer of~\eqref{eq-qot-entropic-primal} is equivalently the minimizer of $D_H(T|K)$ over $\mathcal M_A\cap\mathcal M_B$. Moreover, if a current positive definite matrix has Gibbs form $T_e(F,G)$, then its Bregman projection onto $\mathcal M_A$ has the form $T_e(F^+,G)$, where $F^+$ is chosen so that $\operatorname{Tr}_B T_e(F^+,G)=A$. The projection onto $\mathcal M_B$ is analogous. Thus, when each one-block marginal equation is solved exactly, alternating Bregman projections are equivalent to alternating block maximization of the dual~\eqref{eq-qot-entropic-dual}.
\end{prop}
\begin{proof}
Since $\log K=-C/\epsilon$, the identity
\(\tr(CT)+\epsilon H(T) = \epsilon D_H(T|K)-\epsilon\tr(K)\)
holds, so the primal minimizer is the constrained Bregman projection of $K$ up to an additive constant. For the projection of a positive definite matrix $S$ onto $\mathcal M_A$, the affine set contains the positive definite point $A\otimes\Id_m/m$; the entropy derivative $\log T-\log S$ is singular at the boundary, so the projection lies in the interior of the positive cone. Its Lagrangian first variation is
\(\log T-\log S-\Lambda\otimes\Id_m=0\)
for a Hermitian multiplier $\Lambda$. Hence $T=\exp(\log S+\Lambda\otimes\Id_m)$. If $S=T_e(F,G)$, this is again of the form $T_e(F+\epsilon\Lambda,G)$. The multiplier is fixed by the marginal equation $\operatorname{Tr}_B(T)=A$. The same argument applies to $\mathcal M_B$. Finally, the first-order optimality condition for maximizing~\eqref{eq-qot-entropic-dual} over one block is exactly the corresponding marginal equation, so the Bregman and block-dual views coincide.
\end{proof}
\(\operatorname{Tr}_B T_e(F,G)=A, \qquad \operatorname{Tr}_A T_e(F,G)=B\)
\paragraph{Gurvits scaling and quantum Sinkhorn.}
\begin{equation}\label{eq-qot-symmetric-scaling}
T_s(F,G)
=
\exp\!\left(\frac{Z}{2\epsilon}\right)\exp(-C/\epsilon)\exp\!\left(\frac{Z}{2\epsilon}\right)
=
(U\otimes V)K(U\otimes V),
\qquad
Z=F\otimes\Id_m+\Id_n\otimes G,
\end{equation}
where $U=\exp(F/(2\epsilon))$, $V=\exp(G/(2\epsilon))$ and $K=\exp(-C/\epsilon)$. If $[Z,C]=0$, then $T_s(F,G)=T_e(F,G)$; otherwise this is a Strang-type symmetric surrogate.

Fix a Choi convention and let $\mathcal K:\mathbb H_m\to\mathbb H_n$ be the completely positive map represented by the positive Choi matrix $K$; let $\mathcal K^\star$ be its Hilbert--Schmidt adjoint. Up to the transpose dictated by the chosen Choi convention, the marginal equations for the symmetric coupling take the operator-scaling form

\(U\,\mathcal K(V^2)\,U=A, \qquad V\,\mathcal K^\star(U^2)\,V=B.\)
\begin{equation}\label{eq-qot-gurvits-updates}
\begin{aligned}
R_V&=\mathcal K(V^2),
&
U&\leftarrow
R_V^{-1/2}\bigl(R_V^{1/2} A R_V^{1/2}\bigr)^{1/2}R_V^{-1/2},
\\
S_U&=\mathcal K^\star(U^2),
&
V&\leftarrow
S_U^{-1/2}\bigl(S_U^{1/2} B S_U^{1/2}\bigr)^{1/2}S_U^{-1/2}.
\end{aligned}
\end{equation}
These inverse square roots are well-defined when $K\succ0$ and $U,V,A,B\succ0$. This is Gurvits/operator scaling with prescribed targets.